\documentclass[11pt]{article}

\usepackage[a4paper,margin=22mm]{geometry}
\usepackage{fontspec}
\usepackage{xeCJK}
\usepackage{amsmath,amssymb,mathtools,bm}
\usepackage{booktabs}
\usepackage{enumitem}
\usepackage{xcolor}
\usepackage[hidelinks]{hyperref}
\usepackage{fancyhdr}

\setmainfont{Helvetica Neue}
\setsansfont{Helvetica Neue}
\setmonofont{Menlo}
\setCJKmainfont{Songti SC}
\setCJKsansfont{PingFang SC}
\setCJKmonofont{PingFang SC}

\definecolor{ink}{HTML}{111111}
\definecolor{muted}{HTML}{555555}
\definecolor{linegray}{HTML}{DDDDDD}

\hypersetup{
  pdftitle={Cyrus CS231A Guided Notes},
  pdfauthor={Cyrus Learning Manager},
  pdfsubject={Stanford CS231A guided study notes},
  pdfkeywords={computer vision, SLAM, 3D reconstruction, spatial intelligence}
}

\pagestyle{fancy}
\fancyhf{}
\lhead{\small Cyrus CS231A Guided Notes}
\rhead{\small \thepage}
\renewcommand{\headrulewidth}{0.4pt}

\setlength{\parindent}{0pt}
\setlength{\parskip}{0.55em}
\setlist[itemize]{leftmargin=1.4em,itemsep=0.2em,topsep=0.2em}
\setlist[enumerate]{leftmargin=1.6em,itemsep=0.2em,topsep=0.2em}
\sloppy

\newcommand{\R}{\mathbb{R}}
\newcommand{\K}{\mathbf{K}}
\newcommand{\E}{\mathbf{E}}
\newcommand{\F}{\mathbf{F}}
\newcommand{\T}{\mathbf{T}}
\newcommand{\M}{\mathbf{M}}
\newcommand{\x}{\mathbf{x}}
\newcommand{\p}{\mathbf{p}}
\newcommand{\X}{\mathbf{X}}
\newcommand{\Normal}{\mathcal{N}}
\newcommand{\goodnotes}[1]{\textbf{GoodNotes 任务：}#1}

\begin{document}

\begin{titlepage}
  \centering
  {\sffamily\bfseries\Huge Cyrus CS231A 导学讲义\par}
  \vspace{0.8em}
  {\Large Computer Vision, 3D Reconstruction, SLAM, Spatial Intelligence\par}
  \vspace{1.2em}
  {\large 基于 Stanford CS231A course notes 的中文学习版整理\par}
  \vfill
  \begin{tabular}{ll}
    官方课程 & \href{https://web.stanford.edu/class/cs231a/}{Stanford CS231A} \\
    本地 PDF & \texttt{public/courses/cs231a/course-notes/} \\
    学习方式 & 网页交互 + GoodNotes 推导 + Obsidian 思维图 + Notion 进度证据 \\
  \end{tabular}
  \vfill
  {\small 本文是学习导学笔记，不替代官方讲义。公式使用 LaTeX 源码编写并在 PDF 中渲染。}
\end{titlepage}

\tableofcontents
\newpage

\section*{使用方式}
\addcontentsline{toc}{section}{使用方式}

这份 PDF 不是每日计划，而是自学时随时打开的路线图。每一章对应一份 CS231A course note。你在 iPad Pro 的 GoodNotes 里做两件事：
\begin{enumerate}
  \item 抄核心公式，不要求马上全会，但要知道每个符号代表什么。
  \item 画一张几何直觉图，例如相机、射线、平面、点云、体素、状态估计链。
\end{enumerate}

Notion 负责记录证据：当前章节、GoodNotes 页码、掌握程度、仍不懂的问题。Obsidian 负责把概念连成图：camera model $\rightarrow$ epipolar geometry $\rightarrow$ triangulation $\rightarrow$ SfM $\rightarrow$ SLAM/state estimation。

\section{课程总地图}

\begin{center}
\begin{tabular}{p{0.08\linewidth}p{0.28\linewidth}p{0.55\linewidth}}
\toprule
编号 & 讲义 & 学习目标 \\
\midrule
1 & Camera Models & 从 3D 点到 2D 像素，建立相机内参、外参和标定基础。 \\
2 & Single View Metrology & 用单张图里的消失点、消失线推回 3D 方向和相机信息。 \\
3 & Epipolar Geometry & 两张图之间的点线约束：$\E$、$\F$、八点法和校正。 \\
4 & Stereo Systems and SfM & 三角测量、重投影误差、SfM、歧义和 Bundle Adjustment。 \\
5 & Active and Volumetric Stereo & 主动双目、体素重建、space carving、voxel coloring。 \\
6 & Fitting and Matching & 最小二乘、鲁棒拟合、RANSAC、Hough voting。 \\
7 & Representation Learning & 状态、表示学习、生成式/判别式、自监督表示。 \\
8 & Monocular Depth and Tracking & 单目深度、光度重建损失、特征跟踪、dense descriptors。 \\
9 & Optical and Scene Flow & 光流约束、孔径问题、Lucas-Kanade、scene flow 入口。 \\
10 & Optimal Estimation & Bayes filter、Kalman filter、EKF，连接 SLAM 后端。 \\
\bottomrule
\end{tabular}
\end{center}

\section{Camera Models}

\subsection{从针孔模型开始}

相机模型的第一件事是把 3D 点投到 2D 平面。若 3D 点为 $(x,y,z)$，焦距为 $f$，针孔投影为
\[
  x' = f\frac{x}{z}, \qquad y' = f\frac{y}{z}.
\]
核心直觉：深度 $z$ 在分母里，所以物体越远，投影越小。单个像素并不能唯一确定 3D 点，它只确定一条从相机中心出发的射线。

\subsection{从物理成像到像素坐标}

数字图像不是连续平面，而是像素网格。主点为 $(c_x,c_y)$，像素尺度吸收到 $\alpha,\beta$ 后：
\[
  u = \alpha \frac{x}{z}+c_x,\qquad
  v = \beta \frac{y}{z}+c_y.
\]
相机内参矩阵为
\[
  \K =
  \begin{bmatrix}
    \alpha & 0 & c_x\\
    0 & \beta & c_y\\
    0 & 0 & 1
  \end{bmatrix}.
\]
若考虑 skew，可写成
\[
  \K =
  \begin{bmatrix}
    \alpha & -\alpha \cot \theta & c_x\\
    0 & \beta/\sin\theta & c_y\\
    0 & 0 & 1
  \end{bmatrix}.
\]

\subsection{完整相机投影}

外参描述世界坐标到相机坐标：
\[
  P_c = R P_w + T.
\]
完整投影写成
\[
  s
  \begin{bmatrix}
    u\\v\\1
  \end{bmatrix}
  =
  \K
  \begin{bmatrix}
    R & T
  \end{bmatrix}
  \begin{bmatrix}
    X_w\\Y_w\\Z_w\\1
  \end{bmatrix}.
\]
一句话记忆：外参管相机在哪里，内参管相机怎么成像。

\subsection{相机标定}

已知一组 3D 标定点 $P_i$ 和图像点 $p_i$，要求相机矩阵 $\M$：
\[
  p_i \sim \M P_i.
\]
令 $\M$ 的三行分别为 $m_1,m_2,m_3$，则
\[
  u_i = \frac{m_1P_i}{m_3P_i},\qquad
  v_i = \frac{m_2P_i}{m_3P_i}.
\]
整理成线性约束：
\[
  u_i(m_3P_i)-m_1P_i=0,\qquad
  v_i(m_3P_i)-m_2P_i=0.
\]
因为 $\M$ 是 $3\times4$ 矩阵且整体尺度不重要，所以有 11 个自由度；每个对应点给 2 个方程，至少需要 6 个 3D--2D 对应点。

\subsection{刚体变换、畸变和弱透视}

旋转矩阵满足
\[
  R^TR=I,\qquad \det(R)=1.
\]
齐次坐标下的刚体变换为
\[
  \begin{bmatrix}
    v'\\1
  \end{bmatrix}
  =
  \begin{bmatrix}
    R&t\\0&1
  \end{bmatrix}
  \begin{bmatrix}
    v\\1
  \end{bmatrix}.
\]
弱透视假设物体深度变化相对平均深度很小：
\[
  x' \approx \frac{f}{z_0}x,\qquad y' \approx \frac{f}{z_0}y.
\]
\goodnotes{画出相机中心、图像平面、3D 点、投影点。标出 $z$ 在分母里的位置。}

\section{Single View Metrology}

\subsection{四类 2D 变换}

刚体变换保长度：
\[
  \begin{bmatrix}
    x'\\y'\\1
  \end{bmatrix}
  =
  \begin{bmatrix}
    R&t\\0&1
  \end{bmatrix}
  \begin{bmatrix}
    x\\y\\1
  \end{bmatrix}.
\]
相似变换保角度和长度比例：
\[
  \begin{bmatrix}
    x'\\y'\\1
  \end{bmatrix}
  =
  \begin{bmatrix}
    sR&t\\0&1
  \end{bmatrix}
  \begin{bmatrix}
    x\\y\\1
  \end{bmatrix}.
\]
仿射变换保直线和平行性：
\[
  T(v)=Av+t.
\]
投影变换保直线，但不保平行性：
\[
  x' \sim Hx.
\]

\subsection{直线、交点和无穷远点}

直线用齐次向量表示：
\[
  l =
  \begin{bmatrix}
    a\\b\\c
  \end{bmatrix},\qquad
  l^T
  \begin{bmatrix}
    x\\y\\1
  \end{bmatrix}=0.
\]
两条线的交点：
\[
  x = l\times l'.
\]
平行线在齐次坐标里会交于无穷远点，其第三个坐标为 0。

\subsection{消失点和消失线}

若 3D 中一组平行线的方向为 $d$，在图像中的消失点为
\[
  v = \K d.
\]
若知道 $\K$，可反推出方向：
\[
  d = \frac{\K^{-1}v}{\|\K^{-1}v\|}.
\]
一个平面上的多组平行线会产生多个消失点，它们连成这个平面的消失线或地平线 $\ell_{\text{horiz}}$。平面法向量可由
\[
  n = \K^T\ell_{\text{horiz}}
\]
估计。

\subsection{用消失点估计角度和内参}

两组 3D 方向 $d_1,d_2$ 的夹角满足
\[
  \cos\theta =
  \frac{d_1\cdot d_2}{\|d_1\|\|d_2\|}
  =
  \frac{v_1^T\omega v_2}
  {\sqrt{v_1^T\omega v_1}\sqrt{v_2^T\omega v_2}},
\qquad
  \omega=(\K\K^T)^{-1}.
\]
若两组方向真实垂直：
\[
  v_1^T\omega v_2=0.
\]
房间、走廊、建筑非常适合单视图测量，因为它们包含大量平行和垂直结构。

\goodnotes{画三组互相垂直的边线，分别求三个消失点，再写下 $v_i^T\omega v_j=0$。}

\section{Epipolar Geometry}

\subsection{为什么需要两张图}

单张图的像素点只确定一条 3D 射线。第二张图加入后，第一张图中点 $\p$ 的对应点 $\p'$ 不再需要在整张图里搜索，而只需落在第二张图的一条对极线 $\ell'$ 上。

两个相机中心 $O_1,O_2$ 和 3D 点 $P$ 确定对极平面。该平面与两张图像平面的交线就是对极线。基本约束为
\[
  \p' \in \ell'.
\]

\subsection{Essential Matrix}

在归一化相机坐标里，相机间相对运动为 $R,T$。本质矩阵为
\[
  \E = [T_\times]R.
\]
叉乘矩阵定义为
\[
  a\times b =
  [a_\times]b =
  \begin{bmatrix}
    0&-a_z&a_y\\
    a_z&0&-a_x\\
    -a_y&a_x&0
  \end{bmatrix}
  b.
\]
对极约束：
\[
  \p^T\E\p'=0.
\]
$\E$ 包含两个相机之间的相对旋转和平移，秩为 2，有 5 个自由度。

\subsection{Fundamental Matrix}

真实图像通常是像素坐标，内参不能忽略。基础矩阵 $\F$ 直接约束像素坐标：
\[
  \p^T\F\p'=0.
\]
与 $\E$ 的关系为
\[
  \F = \K^{-T}\E{\K'}^{-1},\qquad
  \E = \K^T\F\K'.
\]
对极线可由
\[
  \ell'=\F^T\p,\qquad \ell=\F\p'
\]
计算。

\subsection{八点法和归一化八点法}

每对匹配点 $\p_i=(u_i,v_i,1)$ 与 $\p'_i=(u'_i,v'_i,1)$ 给出
\[
  \p_i^T\F\p'_i=0.
\]
展开后形成
\[
  Wf=0,
\]
其中 $f$ 是 $\F$ 拉平后的 9 维向量。$\F$ 只差整体尺度，所以最少需要 8 对匹配点。实际先用 SVD 得到 $\hat{\F}$，再强制
\[
  \operatorname{rank}(\F)=2,\qquad \det(\F)=0.
\]
归一化八点法先做
\[
  q_i=T\p_i,\qquad q'_i=T'\p'_i
\]
来改善数值条件，再反归一化：
\[
  \F={T'}^T\F_qT.
\]

\subsection{图像校正}

图像校正寻找两个单应矩阵 $H_1,H_2$：
\[
  \tilde{\p}=H_1\p,\qquad \tilde{\p}'=H_2\p'.
\]
校正后，对应点在左右图中的纵坐标相同：
\[
  \tilde v=\tilde v'.
\]
这把二维搜索变成一维搜索，是双目深度估计的关键。

\section{Stereo Systems and Structure from Motion}

\subsection{三角测量}

若两个相机都看到 3D 点 $P$，它们的图像点为 $\p,\p'$。已知相机矩阵 $\M,\M'$ 时，理想情况下两条反投影射线相交于 $P$。现实中有噪声，所以用重投影误差：
\[
  \min_{\hat P}
  \|\M\hat P-\p\|^2+\|\M'\hat P-\p'\|^2.
\]
线性三角测量可写为
\[
  AP=0,
\]
其中
\[
  A=
  \begin{bmatrix}
    xM_3-M_1\\
    yM_3-M_2\\
    x'M'_3-M'_1\\
    y'M'_3-M'_2
  \end{bmatrix}.
\]

\subsection{非线性精修}

投影要做齐次归一化，因此精修是非线性问题：
\[
  \min_{\hat P}\sum_i\|\M_i\hat P-\p_i\|^2.
\]
Gauss--Newton 在当前估计附近线性化：
\[
  e(\hat P+\delta P)\approx e(\hat P)+J\delta P.
\]
修正量为
\[
  \delta P=-(J^TJ)^{-1}J^Te,
\qquad
  \hat P\leftarrow \hat P+\delta P.
\]

\subsection{Affine Structure from Motion}

仿射或弱透视模型下：
\[
  x_{ij}=A_iX_j+b_i.
\]
对每张图中心化：
\[
  \hat x_{ij}=x_{ij}-\bar x_i.
\]
平移项消失：
\[
  \hat x_{ij}=A_i\hat X_j.
\]
堆成测量矩阵：
\[
  D=MS,
\]
其中 $M$ 是 motion matrix，$S$ 是 structure matrix。理想情况下
\[
  \operatorname{rank}(D)=3.
\]
SVD 分解：
\[
  D\approx U_3\Sigma_3V_3^T.
\]
Tomasi--Kanade 常用
\[
  M=U_3\sqrt{\Sigma_3},\qquad S=\sqrt{\Sigma_3}V_3^T.
\]

\subsection{重建歧义}

分解 $D=MS$ 不是唯一的，因为任意可逆矩阵 $A$ 可插入：
\[
  D=MAA^{-1}S=(MA)(A^{-1}S).
\]
仿射重建保平行关系，但长度、角度和尺度不一定真实。标定相机下，歧义通常缩小到 similarity ambiguity：整体旋转、平移和尺度。

\subsection{Perspective SfM 与 Bundle Adjustment}

从 $\F$ 可构造一组等价投影矩阵：
\[
  \tilde M_1=[I\mid0],\qquad
  \tilde M_2=[-[e_\times]\F\mid e],
\]
其中 $e$ 是极点，满足 $\F e=0$。若相机已标定：
\[
  \E=\K'^T\F\K,\qquad
  \E=[t_\times]R.
\]
从 $\E$ 分解 $R,t$ 会有 4 组候选；正确解让最多三角化点同时位于两个相机前方。

Bundle Adjustment 同时优化相机和 3D 点：
\[
  \min_{\{M_i\},\{X_j\}}
  \sum_{(i,j)\in\mathcal V}
  \|\pi(M_iX_j)-x_{ij}\|^2.
\]
它是 SfM 和 SLAM 后端优化的核心思想。

\section{Active and Volumetric Stereo}

\subsection{主动双目}

传统双目最大难点是对应关系：
\[
  \p\leftrightarrow\p'.
\]
主动双目把一台相机换成投影仪：
\[
  \text{projector}+\text{camera}.
\]
投影仪发出已知点、线、彩色条纹或红外 pattern，相机观察到对应点后即可三角测量：
\[
  P=\operatorname{triangulate}(p,p').
\]
优点是降低匹配难度；代价是投影仪和相机都要标定，且受反光、透明、强光等影响。

\subsection{体素重建}

Volumetric Stereo 从 3D 空间假设开始：先定义工作体积 $V$，划成 voxel grid，然后把每个 voxel 投影到多张图，检查一致性。

Space carving 使用 silhouette。每张 silhouette 反投影成 visual cone，物体在多个视觉锥交集中：
\[
  \text{visual hull}=\bigcap_i \text{visual cone}_i.
\]
若一个 voxel 投影到任一视图的 silhouette 外，就删掉。

Shadow carving 增加自阴影检查，帮助处理 concavity。Voxel coloring 使用颜色一致性：
\[
  \operatorname{color}(x_1)\approx \operatorname{color}(x_2)\approx\cdots\approx \operatorname{color}(x_n).
\]
它要求表面近似 Lambertian，即观察颜色不随视角大幅变化。

\section{Fitting and Matching}

\subsection{普通最小二乘}

拟合直线
\[
  y=mx+b
\]
时，令
\[
  w=
  \begin{bmatrix}
    m\\b
  \end{bmatrix},
\qquad
  E=\|Y-Xw\|^2.
\]
求导得到 normal equation：
\[
  X^TXw=X^TY,
\]
闭式解为
\[
  w=(X^TX)^{-1}X^TY.
\]

\subsection{Total Least Squares}

普通最小二乘只看 $y$ 方向误差，不适合竖直线。使用一般直线：
\[
  ax+by+c=0.
\]
点到线距离：
\[
  d_i=\frac{|ax_i+by_i+c|}{\sqrt{a^2+b^2}}.
\]
加约束
\[
  a^2+b^2=1
\]
后，误差为
\[
  E=\sum_i(ax_i+by_i+c)^2.
\]
最优直线经过数据中心 $(\bar x,\bar y)$：
\[
  c=-a\bar x-b\bar y.
\]
中心化后用 SVD 解，最小奇异值对应方向给出法向量。

\subsection{鲁棒代价与 RANSAC}

平方误差
\[
  C(u_i)=u_i^2
\]
会让外点影响过大。鲁棒代价示例：
\[
  C(u_i,\sigma)=\frac{u_i^2}{\sigma^2+u_i^2}.
\]
RANSAC 反复随机采最小点集，拟合模型，统计 inliers：
\[
  P=\{x_i\mid r(x_i,w)<\delta\}.
\]
估计基础矩阵时最小样本数为 $s=8$；估计单应矩阵时 $s=4$。迭代次数可估计为
\[
  n=\frac{\log(1-p)}
  {\log(1-(1-\epsilon)^s)}.
\]

\subsection{Hough Transform}

Hough voting 把图像点映射到参数空间。普通斜截式参数无界，常用极坐标直线：
\[
  x\cos\theta+y\sin\theta=\rho.
\]
每个图像点在 $(\rho,\theta)$ 空间投票，峰值对应图像中的直线。

\section{Representations and Representation Learning}

\subsection{状态、Markov 性和观测}

状态 $x_t$ 是系统当前足以预测未来的压缩描述。Markov 性：
\[
  P(x_{t+1}\mid x_t,x_{t-1},\dots,x_1)=P(x_{t+1}\mid x_t).
\]
若状态不可直接观测，引入观测 $z_t$：
\[
  P(z_t\mid x_t,x_{t-1},\dots)=P(z_t\mid x_t).
\]
这就是 HMM、POMDP、SLAM 状态估计的基础。

\subsection{生成式与判别式}

生成式模型描述联合分布：
\[
  p(z,x)=p(z\mid x)p(x).
\]
判别式模型直接描述
\[
  p(x\mid z).
\]
视觉任务里，生成式方法像“给定姿态渲染图像再比较”，判别式方法像“神经网络从图像直接回归姿态”。

\subsection{什么是好的表示}

好的表示应该：
\begin{itemize}
  \item compact：压缩但不丢关键因素；
  \item explanatory：足以解释输入变化；
  \item disentangled：不同因素尽量独立；
  \item hierarchical：低级特征组合成高级概念；
  \item useful：让下游推理、预测、决策更容易。
\end{itemize}

\subsection{自编码器和自监督}

Autoencoder 学习瓶颈表示 $z$：
\[
  \hat X=F(X),\qquad L=\|X-\hat X\|.
\]
自监督学习把输入的一部分当标签，例如用图像下半部分预测上半部分。空间智能中的世界模型也在做类似事情：学习一个能预测、压缩和支持决策的 latent state。

\section{Monocular Depth Estimation and Feature Tracking}

\subsection{深度与视差}

校正双目下，视差为
\[
  d=u_L-u_R.
\]
深度公式：
\[
  Z=\frac{fb}{d},
\]
其中 $f$ 是焦距，$b$ 是基线。近处视差大，远处视差小。

单目深度在几何上欠约束，但可以利用透视、尺度、遮挡、光照等统计线索学习。

\subsection{监督、无监督和自监督深度}

监督单目深度直接最小化预测深度与真值深度的误差。无监督双目深度把深度或视差作为中间表示，用右图重建左图：
\[
  \tilde I_l(u,v)=I_r(u-d_l(u,v),v).
\]
典型损失含三部分：
\[
  C_s=\alpha_{ap}(C^l_{ap}+C^r_{ap})
  +\alpha_{ds}(C^l_{ds}+C^r_{ds})
  +\alpha_{lr}(C^l_{lr}+C^r_{lr}).
\]
外观重建损失：
\[
  C^l_{ap}=
  \frac{1}{N}\sum_{i,j}
  \alpha\frac{1-\operatorname{SSIM}(I^l_{ij},\tilde I^l_{ij})}{2}
  +(1-\alpha)\|I^l_{ij}-\tilde I^l_{ij}\|_1.
\]
边缘感知平滑：
\[
  C^l_{ds}=
  \frac{1}{N}\sum_{i,j}
  |\partial_x d^l_{ij}|e^{-\|\partial_x I^l_{ij}\|}
  +
  |\partial_y d^l_{ij}|e^{-\|\partial_y I^l_{ij}\|}.
\]
左右一致性：
\[
  C^l_{lr}=
  \frac{1}{N}\sum_{i,j}
  |d^l_{ij}-d^r_{i,j+d^l_{ij}}|.
\]

自监督单目视频深度把深度和相邻帧位姿一起预测，通过 novel view synthesis 构造光度误差：
\[
  pe(I_a,I_b)=
  \frac{\alpha}{2}(1-\operatorname{SSIM}(I_a,I_b))
  +(1-\alpha)\|I_a-I_b\|_1,
\]
\[
  L_p=\sum_{t'}pe(I_t,I_{t'\rightarrow t}).
\]

\subsection{Feature Tracking 和 Dense Descriptors}

特征跟踪本质是跨帧对应关系问题。传统方法使用手工 keypoint 和 descriptor；现代方法学习每个像素的 dense descriptor：
\[
  f(I,u)\in\R^D.
\]
匹配点希望 descriptor 距离小：
\[
  L_{\text{matches}}=
  \frac{1}{N_{\text{matches}}}
  \sum D(I_a,u_a,I_b,u_b)^2.
\]
非匹配点希望距离大：
\[
  L_{\text{non-matches}}=
  \frac{1}{N_{\text{non-matches}}}
  \sum \max(0,M-D(I_a,u_a,I_b,u_b)^2).
\]
总损失：
\[
  L=L_{\text{matches}}+L_{\text{non-matches}}.
\]

\section{Optical and Scene Flow}

\subsection{Optical Flow 与 Motion Field}

Optical flow 是视频里每个像素的表观 2D 运动：
\[
  \bm u=
  \begin{bmatrix}
    u\\v
  \end{bmatrix}.
\]
Motion field 是 3D 运动向量投影到图像平面的理想速度场。二者不总相同：光照变化可能产生 optical flow，但真实 3D 点未动；无纹理旋转可能有 motion field 但没有可见 optical flow。

\subsection{光流约束方程}

亮度恒定假设：
\[
  I(x,y,t)=I(x+\Delta x,y+\Delta y,t+\Delta t).
\]
小运动假设下，一阶 Taylor 展开：
\[
  I(x+\Delta x,y+\Delta y,t+\Delta t)
  \approx
  I(x,y,t)+I_x\Delta x+I_y\Delta y+I_t\Delta t.
\]
得到光流约束：
\[
  I_xu+I_yv+I_t=0,
\]
或
\[
  -I_t=\nabla I^T\bm u.
\]
一个像素只有一个方程，却有 $u,v$ 两个未知数，所以欠约束，只能得到 normal flow。

\subsection{孔径问题与 Lucas--Kanade}

沿边缘方向的运动无法由局部窗口唯一确定，这叫 aperture problem。Lucas--Kanade 假设邻域内像素共享同一个 flow：
\[
  A\bm u=b.
\]
最小二乘解：
\[
  \bm u_{\text{ls}}=(A^TA)^{-1}A^Tb.
\]
可解性取决于
\[
  A^TA=
  \begin{bmatrix}
    \sum I_x^2 & \sum I_xI_y\\
    \sum I_xI_y & \sum I_y^2
  \end{bmatrix}.
\]
纹理丰富、梯度方向多样的角点区域最适合跟踪。

\section{Optimal Estimation}

\subsection{状态估计目标}

状态估计要递归地估计
\[
  p(x_t\mid z_{1:t},u_{1:t}),
\]
也叫 belief：
\[
  bel(t).
\]
Markov 转移模型：
\[
  p(x_t\mid x_{0:t-1},z_{1:t-1},u_{1:t})
  =
  p(x_t\mid x_{t-1},u_t).
\]
测量模型：
\[
  p(z_t\mid x_{0:t},z_{1:t-1},u_{1:t})
  =
  p(z_t\mid x_t).
\]

\subsection{Bayes Filter}

Prediction：
\[
  \overline{bel}(x_t)
  =
  \int p(x_t\mid x_{t-1},u_t)bel(x_{t-1})\,dx_{t-1}.
\]
Update：
\[
  bel(x_t)=\eta p(z_t\mid x_t)\overline{bel}(x_t),
\]
其中 $\eta$ 是归一化常数。

离散情形可以写成矩阵：
\[
  \widehat{bel}(t)=T(u_t)bel(t-1),
\]
\[
  bel(t)=
  \frac{M(z_t)\widehat{bel}(t)}
  {\mathbf{1}^TM(z_t)\widehat{bel}(t)}.
\]

\subsection{Kalman Filter}

线性高斯系统：
\[
  x_t=A_tx_{t-1}+B_tu_t+w_t,\qquad
  z_t=C_tx_t+v_t.
\]
噪声：
\[
  w_t\sim\Normal(0,Q_t),\qquad v_t\sim\Normal(0,R_t).
\]
预测均值和协方差：
\[
  \mu_{t|t-1}=A_t\mu_{t-1|t-1}+B_tu_t,
\]
\[
  \Sigma_{t|t-1}=A_t\Sigma_{t-1|t-1}A_t^T+Q_t.
\]
更新：
\[
  K_t=
  \Sigma_{t|t-1}C_t^T
  (C_t\Sigma_{t|t-1}C_t^T+R_t)^{-1},
\]
\[
  \mu_{t|t}=\mu_{t|t-1}+K_t(z_t-C_t\mu_{t|t-1}),
\]
\[
  \Sigma_{t|t}=\Sigma_{t|t-1}-K_tC_t\Sigma_{t|t-1}.
\]
这里
\[
  z_t-C_t\mu_{t|t-1}
\]
叫 innovation，表示真实观测和预测观测的差。

\subsection{Extended Kalman Filter}

非线性系统：
\[
  x_t=f(x_{t-1},u_t)+w_t,\qquad
  z_t=g(x_t)+v_t.
\]
EKF 在当前均值附近线性化：
\[
  A_t=
  \left.
  \frac{\partial f(x_t,u_t)}{\partial x_t}
  \right|_{x_t=\mu_{t-1|t-1}},
\qquad
  C_t=
  \left.
  \frac{\partial g(x_t)}{\partial x_t}
  \right|_{x_t=\mu_{t|t-1}}.
\]
然后使用 Kalman filter 的预测和更新形式。SLAM 中许多经典滤波器都可以看成这个框架的变体。

\section{SLAM 与三维重建学习路线}

\begin{enumerate}
  \item 先会相机：$\K,[R\mid T]$，知道像素如何来自 3D 点。
  \item 再会两视图：$\E,\F$，知道点如何约束对极线。
  \item 再会深度：rectification、disparity、triangulation。
  \item 再会鲁棒前端：matching、RANSAC、optical flow。
  \item 再会后端：reprojection error、Gauss--Newton、Bundle Adjustment。
  \item 最后接状态估计：Bayes filter、Kalman filter、EKF、factor graph。
\end{enumerate}

你学习时不需要一天学完。每次只要完成一个闭环：
\[
  \text{公式} \rightarrow \text{几何图} \rightarrow \text{一句话解释} \rightarrow \text{一个失败模式}.
\]

\section*{公式速查}
\addcontentsline{toc}{section}{公式速查}

\begin{align*}
  &\text{Pinhole:} &
  x'&=f\frac{x}{z},&
  y'&=f\frac{y}{z}.\\
  &\text{Camera projection:} &
  s\begin{bmatrix}u\\v\\1\end{bmatrix}
  &=\K[R\mid T]\begin{bmatrix}X\\Y\\Z\\1\end{bmatrix}.\\
  &\text{Vanishing point:} &
  v&=\K d,&
  d&=\frac{\K^{-1}v}{\|\K^{-1}v\|}.\\
  &\text{Epipolar:} &
  \p^T\E\p'&=0,&
  \p^T\F\p'&=0.\\
  &\text{Essential:} &
  \E&=[t_\times]R.\\
  &\text{Fundamental:} &
  \F&=\K^{-T}\E{\K'}^{-1}.\\
  &\text{Stereo depth:} &
  Z&=\frac{fb}{d}.\\
  &\text{Triangulation:} &
  AP&=0.\\
  &\text{Bundle adjustment:} &
  \min_{\{M_i\},\{X_j\}}
  &\sum_{(i,j)\in\mathcal V}
  \|\pi(M_iX_j)-x_{ij}\|^2.\\
  &\text{RANSAC iterations:} &
  n&=\frac{\log(1-p)}{\log(1-(1-\epsilon)^s)}.\\
  &\text{Optical flow:} &
  I_xu+I_yv+I_t&=0.\\
  &\text{Kalman gain:} &
  K_t&=\Sigma_{t|t-1}C_t^T(C_t\Sigma_{t|t-1}C_t^T+R_t)^{-1}.
\end{align*}

\section*{GoodNotes 页码建议}
\addcontentsline{toc}{section}{GoodNotes 页码建议}

\begin{center}
\begin{tabular}{p{0.12\linewidth}p{0.34\linewidth}p{0.44\linewidth}}
\toprule
页码 & 标题 & 必画图 \\
\midrule
001 & Pinhole Camera & 相机中心、图像平面、$z$ 深度。 \\
002 & Camera Matrix & $\K$、$R,T$、世界坐标到像素。 \\
003 & Vanishing Point & 马路两侧线汇聚到消失点。 \\
004 & Epipolar Geometry & 点 $\rightarrow$ 对极线。 \\
005 & Eight-Point and RANSAC & 匹配点、内点、外点。 \\
006 & Triangulation & 两条射线交出 3D 点。 \\
007 & SfM and BA & 相机轨迹、稀疏点云、重投影误差。 \\
008 & Volumetric Stereo & voxel grid 和 visual hull。 \\
009 & Depth and Flow & 视差、光流向量、孔径问题。 \\
010 & Bayes/Kalman/EKF & predict-update 状态估计链。 \\
\bottomrule
\end{tabular}
\end{center}

\end{document}
